La clase anterior terminó con la definición de espectro secuencial, sus dos fuentes principales ---el espectro suspensión $\Sigma^\infty A$ de un espacio punteado y el espectro $E$ asociado a una teoría de cohomología reducida $\widetilde E^{\,*}$--- y la construcción del desplazamiento $sh^d$. Esta clase empieza el trabajo con esa maquinaria: se fijan los tipos especiales de espectros que aparecerán, se define la categoría de espectros secuenciales y sus morfismos, y se introduce la noción de grupo de homotopía de un espectro, que es la que permite decir cuándo dos espectros son «el mismo» sin necesidad de exhibir un isomorfismo.

El objetivo declarado al comienzo de la sesión es mostrar que \emph{la categoría homotópica estable es estable}, es decir, que en ella la suspensión es invertible. La segunda mitad de la clase se dedica al primer paso de esa demostración: construir la suspensión de un espectro y comprobar que coincide con el desplazamiento $sh^1$ en el nivel de los grupos de homotopía.

%%%%%%%%%%%%%%%%%%%%%%%%%%%%%%%%%%%%%%%%%%%%%%%%%%%%%%%%%%%%%%%%%%%%%%%%%%%%%%%
\section{La categoría de espectros secuenciales}
%%%%%%%%%%%%%%%%%%%%%%%%%%%%%%%%%%%%%%%%%%%%%%%%%%%%%%%%%%%%%%%%%%%%%%%%%%%%%%%

Recordemos que un espectro secuencial es una sucesión de espacios punteados $\{X_n\}_{n\ge0}$ junto con aplicaciones punteadas $\sigma_n\colon\Sigma X_n\to X_{n+1}$, sus \emph{morfismos estructurales}; en la versión $\mathbb Z$-secuencial la sucesión se indexa sobre $\mathbb Z$. En la pizarra los morfismos estructurales se escriben $\xi_n$; en estas notas se usa $\sigma_n$, que es la letra empleada en el capítulo anterior, y $\xi_n$ solo cuando se cita la notación literal de la pizarra.

Conviene fijar de entrada una notación que aparecerá repetidamente. Si $X$ es un espacio \emph{no} punteado, se escribe
\[
\Sigma^\infty_+X\coloneqq\Sigma^\infty(X_+),
\]
donde $X_+$ es $X$ con un punto base disjunto. Así, $\Sigma^\infty_+X$ es un espectro secuencial construido a partir de un espacio que no venía punteado, lo que permite aplicar la maquinaria a espacios como las variedades cerradas de la clase anterior.

%%%%%%%%%%%%%%%%%%%%%%%%%%%%%%%%%%%%%%%%%%%%%%%%%%%%%%%%%%%%%%%%%%%%%%%%%%%%%%%
\subsection{Tipos especiales de espectros}
%%%%%%%%%%%%%%%%%%%%%%%%%%%%%%%%%%%%%%%%%%%%%%%%%%%%%%%%%%%%%%%%%%%%%%%%%%%%%%%

Los espectros que aparecen en la práctica tienen propiedades adicionales que conviene nombrar por separado, porque simplifican los cálculos y porque cada modelo de la categoría homotópica estable se construye privilegiando unas u otras.

\begin{defi}\label{defi:omega-espectro}
Un espectro secuencial $X$ es un \emph{$\Omega$-espectro} si, para todo $n$, la aplicación adjunta del morfismo estructural,
\[
\sigma_n^{\perp}\colon X_n\longrightarrow\Omega X_{n+1},
\]
es una equivalencia débil. Es un \emph{$\Omega$-espectro estricto} si $\sigma_n^{\perp}$ es un homeomorfismo para todo $n$.
\end{defi}

Los espectros que provienen de una teoría de cohomología reducida son $\Omega$-espectros: es exactamente el contenido del \myref{lema:omega} del capítulo anterior, donde se vio que las aplicaciones $E_n\to\Omega E_{n+1}$ inducidas por los isomorfismos de suspensión son equivalencias homotópicas.

\begin{defi}\label{defi:espectro-celular}
Un espectro secuencial $X$ es un \emph{espectro celular} si cada $X_n$ es un complejo celular generalizado punteado y cada morfismo estructural $\sigma_n$ es la inclusión de un subcomplejo. Es un \emph{CW-espectro} si los $X_n$ son complejos CW.
\end{defi}

La condición es de naturaleza técnica: garantiza que los espectros se comporten bien respecto de las construcciones homotópicas, del mismo modo que los complejos CW lo hacen entre los espacios. El ejemplo básico es inmediato.

\begin{ejem}\label{ejem:suspension-cw}
Si $A$ es un CW-complejo punteado, entonces $\Sigma^\infty A$ es un CW-espectro.
\end{ejem}

\begin{proof}
Por definición $(\Sigma^\infty A)_n=\Sigma^nA$, que es un CW-complejo punteado, y el morfismo estructural $\sigma_n\colon\Sigma\Sigma^nA\to\Sigma^{n+1}A$ es la identidad una vez identificadas las dos formas de suspender, que es la inclusión de $\Sigma^nA$ en su suspensión. \qedhere
\end{proof}

%%%%%%%%%%%%%%%%%%%%%%%%%%%%%%%%%%%%%%%%%%%%%%%%%%%%%%%%%%%%%%%%%%%%%%%%%%%%%%%
\subsection{La categoría \texorpdfstring{$\Sp_N$}{Sp_N} y sus morfismos}
%%%%%%%%%%%%%%%%%%%%%%%%%%%%%%%%%%%%%%%%%%%%%%%%%%%%%%%%%%%%%%%%%%%%%%%%%%%%%%%

Con los espectros como objetos se forma una categoría, y la decisión importante es qué se toma como morfismo. La elección es la más rígida posible: se pide que la conmutatividad con los morfismos estructurales sea \emph{estricta}, no salvo homotopía.

\begin{defi}\label{defi:spn}
La categoría $\Sp_N$ tiene por objetos los espectros secuenciales. Sus morfismos de $\{E_n\}$ a $\{F_n\}$ son las familias $(f_n\colon E_n\to F_n)_n$ de aplicaciones punteadas para las que el diagrama
\[
\begin{tikzcd}[column sep=large, row sep=large]
\Sigma E_n \arrow[r,"\sigma_n"] \arrow[d,"\Sigma f_n"'] & E_{n+1} \arrow[d,"f_{n+1}"]\\
\Sigma F_n \arrow[r,"\varsigma_n"'] & F_{n+1}
\end{tikzcd}
\]
conmuta para todo $n$, donde $\sigma_n$ y $\varsigma_n$ son los morfismos estructurales de $E$ y de $F$ respectivamente. Es decir,
\[
\hom_{\Sp_N}\bigl(\{E_n\},\{F_n\}\bigr)=
\Bigl\{(f_n)_n\ \Bigm|\ \text{los cuadrados anteriores conmutan para todo }n\Bigr\}.
\]
La composición se realiza nivel a nivel: $(g_n)\circ(f_n)=(g_n\circ f_n)$.
\end{defi}

Que la composición esté bien definida es inmediato: el compuesto de dos cuadrados conmutativos es un cuadrado conmutativo, y la asociatividad y las identidades se heredan nivel a nivel de $\Topstar$. En la pizarra el morfismo estructural del segundo espectro se escribió con un glifo distinto, lo que refleja que son espectros distintos y no que haya dos nociones diferentes.

\begin{defi}\label{defi:equiv-nivel}
Un morfismo $f\colon E\to F$ de espectros secuenciales es una \emph{equivalencia nivel a nivel} si cada $f_n\colon E_n\to F_n$ es una equivalencia débil de espacios.
\end{defi}

Esta noción es la primera candidata a «isomorfismo» en $\Sp_N$, pero resulta demasiado fina para los propósitos del curso: es rígida en cada nivel y no detecta la información estable. La noción correcta se introduce en la \myref{defi:equiv-estable}, una vez que se dispone de los grupos de homotopía.

%%%%%%%%%%%%%%%%%%%%%%%%%%%%%%%%%%%%%%%%%%%%%%%%%%%%%%%%%%%%%%%%%%%%%%%%%%%%%%%
\section{Los grupos de homotopía de un espectro}
%%%%%%%%%%%%%%%%%%%%%%%%%%%%%%%%%%%%%%%%%%%%%%%%%%%%%%%%%%%%%%%%%%%%%%%%%%%%%%%

Un espectro no es un espacio, de modo que no tiene grupos de homotopía en el sentido ordinario. Sin embargo, sí tiene una sucesión de grupos que hace sus veces, y que se define tomando el colímite de los grupos de homotopía de sus niveles. La construcción está diseñada para que, en el caso del espectro suspensión, devuelva los grupos de homotopía estable del espacio de partida.

%%%%%%%%%%%%%%%%%%%%%%%%%%%%%%%%%%%%%%%%%%%%%%%%%%%%%%%%%%%%%%%%%%%%%%%%%%%%%%%
\subsection{Definición de los grupos de homotopía}
%%%%%%%%%%%%%%%%%%%%%%%%%%%%%%%%%%%%%%%%%%%%%%%%%%%%%%%%%%%%%%%%%%%%%%%%%%%%%%%

\begin{defi}\label{defi:pi-espectro}
Sea $X$ un espectro secuencial y $k\in\mathbb Z$. Se define
\[
\pi_k(X)\coloneqq\varinjlim_{\substack{n\\ n+k\ge0}}\pi_{n+k}(X_n),
\]
donde el sistema dirigente tiene por aplicaciones
\[
\pi_{n+k}(X_n)\longrightarrow\pi_{n+k+1}(X_{n+1}),
\qquad
[r]\longmapsto\bigl[(\sigma_n)_*\circ\Sigma(r)\bigr],
\]
es decir, se suspende la clase y se aplica el morfismo inducido por el morfismo estructural.
\end{defi}

La restricción $n+k\ge0$ es la que hace que todos los grupos del sistema estén definidos, y a medida que $n$ crece la condición se satisface automáticamente. El colímite es sobre un sistema dirigente, así que $\pi_k(X)$ está bien definido como conjunto; más adelante se verá que es un grupo abeliano.

El ejemplo que justifica la definición es el del espectro suspensión.

\begin{ejem}\label{ejem:pi-suspension-infinito}
Para todo espacio punteado $A$ y todo $k\in\mathbb Z$,
\[
\pi_k(\Sigma^\infty A)\cong\pi_k^s(A).
\]
\end{ejem}

\begin{proof}
Por definición $(\Sigma^\infty A)_n=\Sigma^nA$, de modo que
\[
\pi_k(\Sigma^\infty A)=\varinjlim_{n}\pi_{n+k}(\Sigma^nA),
\]
y el morfismo del sistema es el isomorfismo de suspensión $\pi_{n+k}(\Sigma^nA)\cong\pi_{n+k+1}(\Sigma^{n+1}A)$, que en este caso es un isomorfismo y no solo una aplicación. El colímite de un sistema cuyas aplicaciones son isomorfismos es el grupo común, que es precisamente la definición de $\pi_k^s(A)$. \qedhere
\end{proof}

%%%%%%%%%%%%%%%%%%%%%%%%%%%%%%%%%%%%%%%%%%%%%%%%%%%%%%%%%%%%%%%%%%%%%%%%%%%%%%%
\subsection{Un paréntesis categórico: adjunciones}
%%%%%%%%%%%%%%%%%%%%%%%%%%%%%%%%%%%%%%%%%%%%%%%%%%%%%%%%%%%%%%%%%%%%%%%%%%%%%%%

Para calcular $\pi_k$ en el caso de un espectro que proviene de una teoría de cohomología conviene recordar cómo se traduce una adjunción en fórmulas. El resultado es elemental, pero se usará repetidamente.

Sean $F\colon\mathcal C\to\mathcal D$ y $G\colon\mathcal D\to\mathcal C$ funtores con $F\dashv G$, y sea $\varepsilon\colon FG\to\id_{\mathcal D}$ la counidad de la adjunción. Para $f\colon X\to GY$ escribimos $f^{\perp}\colon FX\to Y$ su adjunta. Entonces, para $r\colon Z\to X$,
\[
\operatorname{adj}\bigl(f_*(r)\bigr)=\operatorname{adj}(f\circ r)=\varepsilon_Y\circ F(f\circ r),
\qquad
(f^{\perp})_*\bigl(F(r)\bigr)=f^{\perp}\circ F(r)=\varepsilon_Y\circ F(f)\circ F(r),
\]
y como $F$ es un funtor las dos expresiones coinciden. Es decir: \emph{adjuntar y después componer es lo mismo que componer y después adjuntar}, lo que se resume en la conmutatividad del diagrama
\[
\begin{tikzcd}[column sep=large, row sep=large]
\hom_{\mathcal C}(Z,X) \arrow[r,"f_*"] \arrow[d,"F"'] & \hom_{\mathcal C}(Z,GY) \arrow[d,"\operatorname{adj}"]\\
\hom_{\mathcal D}(FZ,FX) \arrow[r,"(f^{\perp})_*"'] & \hom_{\mathcal D}(FZ,Y)
\end{tikzcd}
\]
donde la flecha izquierda es la aplicación inducida por $F$ y la derecha la biyección de la adjunción.

En la situación que nos ocupa, $F=\Sigma$ y $G=\Omega$. Los morfismos estructurales de un espectro dan, por adjunción, aplicaciones
\[
\sigma_n^{\perp}\colon X_n\longrightarrow\Omega X_{n+1},
\]
y el diagrama anterior permite describir las aplicaciones del sistema dirigente de la \myref{defi:pi-espectro} sin mencionar explícitamente el morfismo estructural.

\begin{nota}
   Conviene tener a la vista la forma concreta del diagrama, porque en la pizarra se escribió con las dos diagonales y es donde se ve que las dos descripciones coinciden. Con $\pi_{n+k}(X_n)$ arriba a la izquierda y $\pi_{n+k+1}(X_{n+1})$ arriba a la derecha, las cuatro flechas son
   \[
   (\sigma_n^{\perp})_*\colon\pi_{n+k}(X_n)\to\pi_{n+k}(\Omega X_{n+1}),
   \qquad
   \operatorname{adj}\colon\pi_{n+k}(\Omega X_{n+1})\to\pi_{n+k+1}(X_{n+1}),
   \]
   \[
   \Sigma\colon\pi_{n+k}(X_n)\to\pi_{n+k+1}(\Sigma X_n),
   \qquad
   (\sigma_n)_*\colon\pi_{n+k+1}(\Sigma X_n)\to\pi_{n+k+1}(X_{n+1}),
   \]
   y el compuesto de las dos primeras coincide con el de las dos últimas, que es la aplicación del sistema dirigente. En la pizarra se anotó al lado que $(\sigma_n)_*$ es, en efecto, la aplicación inducida por $f$, con lo que las dos descripciones quedan identificadas.
\end{nota}

%%%%%%%%%%%%%%%%%%%%%%%%%%%%%%%%%%%%%%%%%%%%%%%%%%%%%%%%%%%%%%%%%%%%%%%%%%%%%%%
\subsection{El espectro asociado a una teoría de cohomología}
%%%%%%%%%%%%%%%%%%%%%%%%%%%%%%%%%%%%%%%%%%%%%%%%%%%%%%%%%%%%%%%%%%%%%%%%%%%%%%%

Con el paréntesis anterior se calculan los grupos de homotopía de un espectro que proviene de una teoría de cohomología reducida, y el resultado tiene una interpretación limpia.

\begin{prop}\label{prop:pi-coeficientes}
Sea $E$ el espectro asociado a una teoría de cohomología reducida $\widetilde E^{\,*}$, de modo que $\widetilde E^{\,n}(X)\cong\set*{X,E_n}$ naturalmente. Entonces, para todo $k\in\mathbb Z$,
\[
\pi_k(E)\cong\widetilde E^{\,-k}(S^0)\cong\widetilde E^{\,0}(S^k).
\]
\end{prop}

\begin{proof}
Por la \myref{defi:pi-espectro} y la representabilidad,
\[
\pi_k(E)=\varinjlim_n\pi_{n+k}(E_n)=\varinjlim_n\bigl[S^{n+k},E_n\bigr]_*\cong\varinjlim_n\widetilde E^{\,n}(S^{n+k}).
\]
Ahora bien, los morfismos estructurales $E_n\to\Omega E_{n+1}$ son equivalencias, de modo que el isomorfismo de suspensión
\[
\widetilde E^{\,n}(S^{n+k})\xrightarrow{\ \cong\ }\widetilde E^{\,n+1}(S^{n+k+1})
\]
es un isomorfismo para todo $n$, y el colímite de un sistema de isomorfismos es el grupo común $\widetilde E^{\,n}(S^{n+k})$ para cualquier $n$. Tomando $n=0$ se obtiene $\widetilde E^{\,0}(S^k)$, y tomando $n=-k$ se obtiene $\widetilde E^{\,-k}(S^0)$. \qedhere
\end{proof}

Los grupos $\pi_k(E)$ son, por tanto, los \emph{coeficientes} de la teoría de cohomología reducida $\widetilde E^{\,*}$, y el resultado dice que el espectro los codifica por completo. Dos ejemplos lo ilustran.

\begin{ejem}\label{ejem:pi-ha-ku}
Para el espectro de Eilenberg--Mac Lane $HA$ y para la $K$-teoría compleja $KU$,
\[
\pi_k(HA)\cong
\begin{cases}
A, & k=0,\\
0, & k\neq0,
\end{cases}
\qquad\qquad
\pi_k(KU)\cong
\begin{cases}
\mathbb Z, & k\ \text{par},\\
0, & k\ \text{impar}.
\end{cases}
\]
\end{ejem}

En el primer caso la teoría es la cohomología singular con coeficientes en $A$, cuyos coeficientes se concentran en grado cero; en el segundo, los coeficientes de la $K$-teoría compleja son $\mathbb Z$ en los grados pares, que es la periodicidad de Bott.

\begin{nota}
   En la pizarra el colímite de $\pi_k(E)$ se escribió con subíndice $k$, mientras que la variable del término es $n$. Es un desliz de notación: el subíndice debe ser $n$, que es el índice que recorre los niveles del espectro, y $k$ es el grado fijo del grupo que se calcula. Se conserva aquí la versión correcta.
\end{nota}

%%%%%%%%%%%%%%%%%%%%%%%%%%%%%%%%%%%%%%%%%%%%%%%%%%%%%%%%%%%%%%%%%%%%%%%%%%%%%%%
\subsection{Equivalencia estable}
%%%%%%%%%%%%%%%%%%%%%%%%%%%%%%%%%%%%%%%%%%%%%%%%%%%%%%%%%%%%%%%%%%%%%%%%%%%%%%%

Un morfismo de espectros induce aplicaciones entre los grupos de homotopía, y eso permite definir la noción correcta de «el mismo espectro».

\begin{prop}\label{prop:f-induce}
Sea $f\colon X\to Y$ un morfismo de espectros secuenciales. Entonces $f$ induce, para todo $k\in\mathbb Z$, un homomorfismo
\[
f_*\colon\pi_k(X)\longrightarrow\pi_k(Y).
\]
\end{prop}

\begin{proof}
Basta comprobar que la familia $(f_n)$ es compatible con los sistemas dirigentes que definen $\pi_k(X)$ y $\pi_k(Y)$. La conmutatividad estricta del diagrama de la \myref{defi:spn} implica que el cuadrado
\[
\begin{tikzcd}[column sep=large, row sep=large]
\pi_{n+k}(X_n) \arrow[r] \arrow[d,"(f_n)_*"'] & \pi_{n+k+1}(X_{n+1}) \arrow[d,"(f_{n+1})_*"]\\
\pi_{n+k}(Y_n) \arrow[r] & \pi_{n+k+1}(Y_{n+1})
\end{tikzcd}
\]
conmuta, pues ambos caminos consisten en suspender y aplicar el morfismo estructural correspondiente, y $f$ conmuta con ellos. Una familia de aplicaciones compatible con dos sistemas dirigentes induce una aplicación entre sus colímites, que es $f_*$. \qedhere
\end{proof}

\begin{defi}\label{defi:equiv-estable}
Un morfismo $f\colon X\to Y$ de espectros secuenciales es una \emph{equivalencia estable} si induce un isomorfismo $f_*\colon\pi_k(X)\to\pi_k(Y)$ para todo $k\in\mathbb Z$. Dos espectros $X$ e $Y$ son \emph{establemente equivalentes} si existe un zigzag de equivalencias estables que los conecta,
\[
X\xrightarrow{\ \simeq\ }X_1\xleftarrow{\ \simeq\ }X_2\xrightarrow{\ \simeq\ }\cdots\xleftarrow{\ \simeq\ }Y .
\]
\end{defi}

La necesidad del zigzag, y no de una sola equivalencia, es el punto delicado de la definición, y el ejemplo de la \myref{subsec:zigzag-necesario} muestra que no es una precaución gratuita.

\begin{ejer}\label{ejer:pi-abeliano}
Demostrar que $\pi_k(X)$ es un grupo abeliano para todo espectro $X$ y todo $k$, y que las aplicaciones $f_*$ de la \myref{prop:f-induce} son homomorfismos de grupos.
\end{ejer}

\begin{sol}
Cada $\pi_{n+k}(X_n)$ es un grupo abeliano, porque $n+k\ge1$ para $n$ suficientemente grande y los grupos de homotopía de un espacio punteado son abelianos en grado positivo. Las aplicaciones del sistema dirigente son homomorfismos de grupos, de modo que $\pi_k(X)$ es el colímite de un sistema dirigente de grupos abelianos y homomorfismos, que es de nuevo un grupo abeliano. Para la segunda afirmación, $f_*$ está inducida por los homomorfismos $(f_n)_*$, y el paso al colímite de homomorfismos es un homomorfismo.
\end{sol}

%%%%%%%%%%%%%%%%%%%%%%%%%%%%%%%%%%%%%%%%%%%%%%%%%%%%%%%%%%%%%%%%%%%%%%%%%%%%%%%
\subsection{El ejemplo de la truncación}
%%%%%%%%%%%%%%%%%%%%%%%%%%%%%%%%%%%%%%%%%%%%%%%%%%%%%%%%%%%%%%%%%%%%%%%%%%%%%%%

El primer ejemplo de equivalencia estable no trivial es el de recortar un espectro $\mathbb Z$-secuencial por debajo de cero. El interés está en que el morfismo va en la dirección contraria a la que uno esperaría, y en que la demostración usa de manera esencial que los grupos de homotopía son colímites.

\begin{ejem}\label{ejem:truncacion}
Sea $X$ un espectro $\mathbb Z$-secuencial. Se define su \emph{truncación} $X_{\ge0}$ por
\[
(X_{\ge0})_n=
\begin{cases}
\ast, & n<0,\\
X_n, & n\ge0,
\end{cases}
\]
con morfismos estructurales los mismos $\sigma_n$ de $X$ para $n\ge0$, la inclusión del punto base $\Sigma\ast\to X_0$ para $n=-1$, y $\Sigma\ast\xrightarrow{1}\ast$ para $n<-1$. Entonces el morfismo
\[
f\colon X_{\ge0}\longrightarrow X,
\qquad
f_n=
\begin{cases}
\text{la inclusión del punto base}, & n<0,\\
1_{X_n}, & n\ge0,
\end{cases}
\]
es una equivalencia estable.
\end{ejem}

\begin{proof}
Hay que ver que $f_*\colon\pi_k(X_{\ge0})\to\pi_k(X)$ es un isomorfismo para todo $k$. Los dos grupos son colímites sobre el mismo sistema de índices,
\[
\pi_k(X_{\ge0})=\varinjlim_{\substack{n\\ n+k\ge0}}\pi_{n+k}\bigl((X_{\ge0})_n\bigr),
\qquad
\pi_k(X)=\varinjlim_{\substack{n\\ n+k\ge0}}\pi_{n+k}(X_n),
\]
y la aplicación entre ellos está inducida por la familia $(f_n)_*$. El punto clave es que los niveles $n<0$ no aportan nada al colímite: como $\pi_{n+k}\bigl((X_{\ge0})_n\bigr)=\pi_{n+k}(\ast)=0$ para $n<0$, se puede restringir el sistema a los índices $n\ge0$ sin cambiar el colímite, porque esos índices forman un subconjunto \emph{cofinal} del sistema de índices. Se obtiene así el cuadrado
\[
\begin{tikzcd}[column sep=huge, row sep=large]
\varinjlim_{\substack{n\\ n+k\ge0}}\pi_{n+k}\bigl((X_{\ge0})_n\bigr) \arrow[r,"(f_n)_*"] \arrow[d,"\cong"'] & \varinjlim_{\substack{n\\ n+k\ge0}}\pi_{n+k}(X_n) \arrow[d,"\cong"]\\
\varinjlim_{\substack{n\ge0\\ n+k\ge0}}\pi_{n+k}\bigl((X_{\ge0})_n\bigr) \arrow[r,"\cong"'] & \varinjlim_{\substack{n\ge0\\ n+k\ge0}}\pi_{n+k}(X_n)
\end{tikzcd}
\]
cuyas flechas verticales son isomorfismos por cofinalidad, y cuya flecha inferior es un isomorfismo porque para $n\ge0$ se tiene $(X_{\ge0})_n=X_n$ y $f_n=1_{X_n}$. Por tanto $f_*$ es un isomorfismo. \qedhere
\end{proof}

Vale la pena subrayar el mecanismo: la truncación \emph{añade} niveles, los que están por debajo de cero, y esos niveles son triviales. El argumento no funcionaría si los niveles añadidos aportaran algo al colímite.

%%%%%%%%%%%%%%%%%%%%%%%%%%%%%%%%%%%%%%%%%%%%%%%%%%%%%%%%%%%%%%%%%%%%%%%%%%%%%%%
\subsection{Por qué el zigzag es necesario}\label{subsec:zigzag-necesario}
%%%%%%%%%%%%%%%%%%%%%%%%%%%%%%%%%%%%%%%%%%%%%%%%%%%%%%%%%%%%%%%%%%%%%%%%%%%%%%%

La \myref{defi:equiv-estable} admite zigzags, y no solo equivalencias, porque en general no existe una equivalencia estable en la dirección que uno querría. El ejemplo más simple es el espectro esfera desplazado.

\begin{ejem}\label{ejem:zigzag-necesario}
Sea $X=sh^1\mathbb S$, el espectro esfera desplazado un lugar. Sus niveles son
\[
X=(\ldots,\ \ast,\ \underset{-1}{S^0},\ \underset{0}{S^1},\ \underset{1}{S^2},\ \ldots),
\]
donde el rótulo bajo cada entrada indica el nivel, y su truncación es
\[
X_{\ge0}=(\ldots,\ \ast,\ \underset{-1}{\ast},\ \underset{0}{S^1},\ \underset{1}{S^2},\ \ldots).
\]
No existe ninguna equivalencia estable $X\to X_{\ge0}$.
\end{ejem}

\begin{proof}
Supongamos que existiera una equivalencia estable $h\colon X\to X_{\ge0}$. Como $h$ es un morfismo de espectros, conmuta con los morfismos estructurales, y en particular el cuadrado
\[
\begin{tikzcd}[column sep=large, row sep=large]
\Sigma S^0 \arrow[r,"1"] \arrow[d,"\Sigma h_{-1}"'] & S^1 \arrow[d,"h_0"]\\
\Sigma\ast \arrow[r,"c"'] & S^1
\end{tikzcd}
\]
conmuta, donde la flecha inferior es la inclusión del punto base. Ahora bien, $\Sigma\ast\cong\ast$, de modo que el camino inferior es constante, y la conmutatividad obliga a que $h_0$ sea constante. Repitiendo el argumento nivel a nivel, $h_1$ es constante, y por inducción $h_n$ es constante para todo $n\ge0$. Por tanto $h$ induce la aplicación nula en todos los grupos de homotopía, y no puede ser una equivalencia estable.

La conclusión no es que $X$ y $X_{\ge0}$ sean distintos: de hecho, sus grupos de homotopía coinciden, pues
\[
\pi_1(X)=\varinjlim_n\pi_{n+1}(X_n)=\varinjlim_n\pi_{n+1}(\Sigma^{n+1}S^0)\cong\mathbb Z\neq0,
\]
y lo mismo ocurre con $X_{\ge0}$. Lo que falla es la existencia de un morfismo que realice el isomorfismo. Por eso la definición de equivalencia estable para espectros, y no solo para morfismos, exige un zigzag. \qedhere
\end{proof}

\begin{nota}
   En la pizarra se anotó al margen, junto a este ejemplo, la forma de calcular los grupos de homotopía de un espectro desplazado:
   \[
   \pi_k(X[d])\cong\pi_{k-d}(X),
   \]
   que se obtiene directamente de la definición del desplazamiento, $X[d]_n=X_{n+d}$, y de la definición de $\pi_k$ como colímite. Es el análogo, para espectros, de la relación $\pi_{k}(\Sigma^dA)\cong\pi_{k-d}(A)$ para espacios.
\end{nota}

%%%%%%%%%%%%%%%%%%%%%%%%%%%%%%%%%%%%%%%%%%%%%%%%%%%%%%%%%%%%%%%%%%%%%%%%%%%%%%%
\subsection{La categoría homotópica estable}
%%%%%%%%%%%%%%%%%%%%%%%%%%%%%%%%%%%%%%%%%%%%%%%%%%%%%%%%%%%%%%%%%%%%%%%%%%%%%%%

Con las equivalencias estables se puede ya construir la categoría que interesa.

\begin{defi}\label{defi:categoria-estable}
La \emph{categoría homotópica estable} es la localización de $\Sp_N$ respecto de las equivalencias estables,
\[
\Sp_N\bigl[\text{equiv.\ estables}^{-1}\bigr],
\]
es decir, la categoría que se obtiene al invertir formalmente todos los morfismos que son equivalencias estables.
\end{defi}

Esta categoría es un modelo de la categoría homotópica estable, y cumple todas las propiedades que se enumeraron en la clase 1. Su ventaja es que no hay que construir nada nuevo: los objetos siguen siendo espectros secuenciales y los morfismos son clases de morfismos. Su inconveniente es el mismo que el de cualquier localización: los morfismos pasan a ser zigzags, y trabajar con clases de zigzags es incómodo. Eso es lo que motiva buscar otros modelos, en los que la inversión sea menos violenta.

\begin{nota}
   En la pizarra se señaló, entre paréntesis, que esta categoría es localmente pequeña, es decir, que los morfismos entre dos objetos forman un conjunto y no una clase propia. Es cierto, pero el argumento no se dio en clase; se cita aquí como dato. También se dejó como ejercicio la comprobación de que el modelo con espectros $\mathbb N$-secuenciales es equivalente al que usa espectros $\mathbb Z$-secuenciales,
   \[
   \Sp_N\bigl[\text{equiv.\ estables}^{-1}\bigr]\simeq\Sp_{\mathbb Z}\bigl[\text{equiv.\ estables}^{-1}\bigr],
   \]
   lo que permite pasar de un convenio de índices al otro según convenga.
\end{nota}

%%%%%%%%%%%%%%%%%%%%%%%%%%%%%%%%%%%%%%%%%%%%%%%%%%%%%%%%%%%%%%%%%%%%%%%%%%%%%%%
\section{La suspensión de espectros}
%%%%%%%%%%%%%%%%%%%%%%%%%%%%%%%%%%%%%%%%%%%%%%%%%%%%%%%%%%%%%%%%%%%%%%%%%%%%%%%

El nombre de «categoría homotópica \emph{estable}» proviene de que en ella la suspensión debe ser invertible. La suspensión está definida sobre espacios, y la primera tarea es extenderla a espectros de modo que sea compatible con la estructura ya construida. La extensión natural no es exactamente el desplazamiento $sh^1$, pero se le parece lo suficiente como para que la diferencia se pueda controlar.

%%%%%%%%%%%%%%%%%%%%%%%%%%%%%%%%%%%%%%%%%%%%%%%%%%%%%%%%%%%%%%%%%%%%%%%%%%%%%%%
\subsection{El funtor de suspensión}
%%%%%%%%%%%%%%%%%%%%%%%%%%%%%%%%%%%%%%%%%%%%%%%%%%%%%%%%%%%%%%%%%%%%%%%%%%%%%%%

\begin{defi}\label{defi:suspension-espectro}
Sea $X$ un espectro secuencial. Su \emph{suspensión} $\Sigma X$ es el espectro secuencial cuyos niveles son
\[
(\Sigma X)_n\coloneqq S^1\wedge X_n
\]
y cuyos morfismos estructurales son
\[
\bigl(S^1\wedge X_n\bigr)\wedge S^1\xrightarrow{\ 1\wedge\sigma_n\ }S^1\wedge X_{n+1},
\]
donde $\sigma_n\colon X_n\wedge S^1\to X_{n+1}$ es el morfismo estructural de $X$, visto como aplicación definida sobre $X_n\wedge S^1$ mediante la identificación $\Sigma X_n=X_n\wedge S^1$.
\end{defi}

El detalle del lado por el que se hace el producto smash es el que hace funcionar la construcción. Los morfismos estructurales de un espectro están definidos sobre $\Sigma X_n=X_n\wedge S^1$, con la esfera \emph{a la derecha}; la suspensión del espectro, en cambio, antepone la esfera, $S^1\wedge X_n$. Esta asimetría es deliberada: es la que permite que el morfismo estructural de $\Sigma X$ se escriba como $1\wedge\sigma_n$ sin necesidad de reordenar factores, y es también la razón de que $\Sigma X$ no sea literalmente $sh^1X$ sino solo isomorfo a él en la categoría homotópica.

Que $\Sigma$ así definido sea un funtor es inmediato: si $f\colon X\to Y$ es un morfismo de espectros, entonces $\Sigma f$ se define nivel a nivel por $\Sigma f_n=1\wedge f_n$, y la conmutatividad con los morfismos estructurales se hereda de la de $f$ y de la funtorialidad del producto smash.

\begin{nota}
   En la pizarra el objetivo de esta sección se resumió con un diagrama en el que aparecían las dos suspensiones y la relación entre ellas:
   \[
   \begin{tikzcd}[column sep=huge, row sep=large]
   \Ho(\Topstar) \arrow[r] \arrow[d,"\Sigma"'] & \Sp_N\bigl[\text{equiv.\ estables}^{-1}\bigr] \arrow[d,"sh^1"]\\
   \Ho(\Topstar) \arrow[r] & \Sp_N\bigl[\text{equiv.\ estables}^{-1}\bigr]
   \end{tikzcd}
   \]
   junto a la anotación de que $\Sigma\neq sh^1$ pero $\Sigma\simeq sh^1$. Es decir: los dos funtores no son iguales sobre los espectros, pero sí son naturalmente isomorfos en la categoría homotópica estable. El resto de la clase se dedica a verificar el primer paso de esa identificación.
\end{nota}

%%%%%%%%%%%%%%%%%%%%%%%%%%%%%%%%%%%%%%%%%%%%%%%%%%%%%%%%%%%%%%%%%%%%%%%%%%%%%%%
\subsection{La suspensión extiende la de espacios}
%%%%%%%%%%%%%%%%%%%%%%%%%%%%%%%%%%%%%%%%%%%%%%%%%%%%%%%%%%%%%%%%%%%%%%%%%%%%%%%

Antes de comparar $\Sigma$ con $sh^1$ conviene comprobar que la definición no es arbitraria: sobre los espectros que provienen de espacios, la nueva suspensión coincide con la suspensión de espacios.

\begin{prop}\label{prop:suspension-extiende}
Para todo espacio punteado $A$,
\[
\Sigma\bigl(\Sigma^\infty A\bigr)=\Sigma^\infty(\Sigma A).
\]
\end{prop}

\begin{proof}
Calculamos los niveles del espectro de la izquierda:
\[
\bigl(\Sigma(\Sigma^\infty A)\bigr)_n=S^1\wedge(\Sigma^\infty A)_n=S^1\wedge\Sigma^nA=S^1\wedge A\wedge(S^1)^{n},
\]
y los del espectro de la derecha:
\[
\bigl(\Sigma^\infty(\Sigma A)\bigr)_n=\Sigma^n(\Sigma A)=S^1\wedge A\wedge(S^1)^{n}.
\]
Las dos expresiones coinciden, y los morfismos estructurales también: en ambos casos el morfismo es $1\wedge1$ una vez identificadas las suspensiones. Es decir, no solo hay un isomorfismo: los dos espectros son iguales. \qedhere
\end{proof}

Es esta igualdad, y no solo un isomorfismo, la que hace que la suspensión de espectros sea una extensión razonable de la de espacios: el diagrama
\[
\begin{tikzcd}[column sep=huge, row sep=large]
A \arrow[r] \arrow[d,"\Sigma"'] & \Sigma^\infty A \arrow[d,"\Sigma"]\\
\Sigma A \arrow[r] & \Sigma^\infty(\Sigma A)=\Sigma(\Sigma^\infty A)
\end{tikzcd}
\]
conmuta estrictamente.

%%%%%%%%%%%%%%%%%%%%%%%%%%%%%%%%%%%%%%%%%%%%%%%%%%%%%%%%%%%%%%%%%%%%%%%%%%%%%%%
\subsection{El cálculo de \texorpdfstring{$\pi_{k+1}(\Sigma X)$}{pi(k+1) de Sigma X}}
%%%%%%%%%%%%%%%%%%%%%%%%%%%%%%%%%%%%%%%%%%%%%%%%%%%%%%%%%%%%%%%%%%%%%%%%%%%%%%%

El primer paso hacia la identificación $\Sigma\simeq sh^1$ es comprobar que los dos funtores producen los mismos grupos de homotopía.

\begin{prop}\label{prop:pi-suspension}
Para todo espectro secuencial $X$ y todo $k\in\mathbb Z$,
\[
\pi_{k+1}(\Sigma X)\cong\pi_k(X).
\]
\end{prop}

\begin{proof}
Por la \myref{defi:pi-espectro}, el grupo $\pi_k(X)$ es el colímite del sistema
\[
\pi_{n+k}(X_n)\xrightarrow{\ \Sigma\ }\pi_{n+k+1}(X_n\wedge S^1)\xrightarrow{\ (\sigma_n)_*\ }\pi_{n+k+1}(X_{n+1}),
\]
y análogamente $\pi_{k+1}(\Sigma X)$ lo es del sistema
\[
\pi_{n+k+1}(S^1\wedge X_n)\xrightarrow{\ \Sigma\ }\pi_{n+k+2}\bigl((S^1\wedge X_n)\wedge S^1\bigr)\xrightarrow{\ (1\wedge\sigma_n)_*\ }\pi_{n+k+2}(S^1\wedge X_{n+1}).
\]
El argumento consiste en comparar los dos sistemas, y tiene tres piezas.

\emph{Primera pieza: paso a los índices pares.} El sistema de índices admite subconjuntos cofinales, y por tanto se puede restringir el colímite a los índices pares $n=2m$ sin alterarlo. La razón de hacerlo es que, como se verá en la tercera pieza, la comparación de los dos sistemas introduce una permutación de factores cuyo signo es $(-1)$; al recorrer dos niveles de golpe, la permutación se aplica dos veces y el signo desaparece. Esta es la observación que se subrayó en clase: la flexibilidad de tomar cualquier subsucesión cofinal es lo que permite arreglar el signo, y es una ventaja que los espectros tienen sobre los espacios.

\emph{Segunda pieza: la permutación.} Para cada $n$ hay una equivalencia
\[
f_n\colon X_n\wedge S^1\longrightarrow S^1\wedge X_n
\]
que intercambia los dos factores, la trenza del producto smash. Estas equivalencias son naturales y compatibles con los morfismos estructurales, de modo que inducen isomorfismos entre los términos de los dos sistemas y, por tanto, entre sus colímites.

\emph{Tercera pieza: el signo.} Falta comprobar que los isomorfismos inducidos por las $f_n$ conmutan con las aplicaciones de los dos sistemas dirigentes. La comprobación es explícita. Como las dos composiciones empiezan por el mismo morfismo, basta comparar el resto, y para ello se escriben las aplicaciones sobre las coordenadas del producto de tres esferas: un elemento $[x,z_1,z_2,z_3]$ se envía por un camino a
\[
\bigl[\sigma_{2m+1}\bigl(\sigma_{2m}(x,z_1),z_2\bigr),\ z_3\bigr],
\]
y por el otro a
\[
\bigl[z_1,\ \sigma_{2m+1}\bigl(\sigma_{2m}(x,z_2),z_3\bigr)\bigr],
\]
de modo que las dos composiciones difieren por la permutación cíclica de los tres factores de $S^1$. Esa permutación es un homeomorfismo de $S^1\wedge S^1\wedge S^1$ cuyo grado es $+1$ cuando se aplica dos veces seguidas ---o, equivalentemente, cuyo signo es $(-1)$ cuando se aplica una sola vez---, y de ahí la necesidad de haber pasado a los índices pares en la primera pieza. Con eso los dos sistemas dirigentes son isomorfos y sus colímites coinciden. \qedhere
\end{proof}

\begin{nota}
   En la pizarra el cálculo se acompañó de un diagrama grande en el que aparecían las dos filas de grupos, una con $S^1$ antepuesta y otra con $S^1$ pospuesta, y las flechas verticales inducidas por las permutaciones $f_n$. Se reprodujo también, en rojo, el desarrollo de las dos composiciones sobre las coordenadas, que es la tercera pieza de la demostración anterior. La conclusión escrita fue la identificación $\pi_{k+1}(\Sigma X)\cong\pi_k(X)$, que es la que se acaba de probar.
\end{nota}

\begin{nota}
   Conviene no confundir lo que se ha probado con lo que falta. La \myref{prop:pi-suspension} dice que $\Sigma X$ y $sh^1X$ tienen los mismos grupos de homotopía, pero \emph{no} exhibe un morfismo entre ellos, y menos aún una equivalencia estable. La construcción del morfismo es un paso posterior, y en clase se anunció que consta de tres etapas: primero la comparación de los grupos de homotopía ---que es la de esta clase---, después la noción de \emph{morfismo débil}, que permite construir un zigzag de morfismos a partir de una sola aplicación definida salvo homotopía, y finalmente la construcción del morfismo mismo. Con eso quedará establecida la identificación $\Sigma\simeq sh^1$ en la categoría homotópica estable.
\end{nota}

El programa de las próximas sesiones queda así trazado. La \myref{defi:categoria-estable} da la categoría, y el trabajo de esta clase ha mostrado que en ella la suspensión se comporta como el desplazamiento de índices, al menos en el nivel de los grupos de homotopía. Lo que falta es construir, para cada espectro $X$, espectros equivalentes a $X$ y a $\Sigma X$ cuyos grupos de homotopía estén dados por los grupos abelianos correspondientes, y después una equivalencia estable entre ellos. Ese es el punto de partida de la clase siguiente.
